\chapter{Introduction and Motivation}
\label{sec:causality-intro}
% ===================================================================

Physics is, at its core, a theory of cause and effect: given the present state
of a system and the laws governing its evolution, what future states are
possible, and with what probabilities?  This paper argues that
\emph{graph-structured lambda theories} (GSLTs) --- the natural categorical
setting for process calculi and rewriting systems --- provide an unusually
transparent foundation for this question.

The central thesis is that a good deal of the apparatus of classical and
quantum mechanics --- least action, path sums, conservation laws, time
symmetry --- is not an additional assumption layered on top of a dynamical
theory, but arises from the combinatorial structure of a GSLT once that
structure is equipped with strikingly little further data.  The data
amounts to:

\begin{itemize}[itemsep=2pt]
  \item a \emph{decoration}: one map assigning to each rewrite step a value
        in a semiring, indexed by a context-decorated Hennessy--Milner
        logic.  Choosing the semiring chooses the theory --- Boolean gives
        bare nondeterminism, tropical gives cost and least action,
        the non-negative reals give stochastic dynamics, the complex
        numbers give amplitudes;
  \item a \emph{conversion structure} on the kinds of resource, whose
        symmetry --- path-independence of conversion --- is what forces an
        energy function into existence at all;
  \item an \emph{authority discipline} saying not merely what may be spent
        but who may spend it.
\end{itemize}

An earlier presentation of this material carried two maps rather than one,
called them a Lagrangian and a Hamiltonian, and claimed that everything
else followed by construction.  The first was an unnecessary duplication;
the second was a misnaming; and the third was an overstatement.  What
follows by construction is less than was claimed and more interesting than
was claimed, because the parts that do \emph{not} follow by construction
turn out to be exactly the parts that carry content: conservation holds
only under a condition, that condition has a cohomological measure, and
the identification of the conserved charge with the generator of time
evolution remains a conjecture.  We have tried to be explicit about which
is which, and the table below is annotated accordingly.

\section{A First Glimpse: The Rosetta Stone}

The following table previews the correspondence developed in this part.
The third column records the \emph{status} of each row, and the reader is
asked to take it seriously: \textsc{con} marks a construction, true by
definition of the objects involved; \textsc{thm} marks a theorem with a
hypothesis that can fail; \textsc{cnj} marks a conjecture we have not
proved.  A fully annotated version appears in
Section~\ref{sec:synthesis}.

\begin{center}
\renewcommand{\arraystretch}{1.35}
\begin{tabular}{>{\itshape}p{4.3cm} p{5.9cm} c}
\toprule
\normalfont\textbf{GSLT} & \textbf{Physics} & \textbf{Status} \\
\midrule
Term                          & State / initial condition & \textsc{con} \\
Rewrite rule                  & Law of motion & \textsc{con} \\
Rewrite path                  & Trajectory & \textsc{con} \\
Synchronization tree          & Causal graph & \textsc{con} \\
Minimum path length           & Proper time & \textsc{con} \\
Equations $\eqs$              & Gauge equivalence & \textsc{con} \\
Reversible envelope $\GSLT^\dagger$ & Time-symmetric theory & \textsc{con} \\
Decoration, tropical          & Action functional; least action & \textsc{thm} \\
Decoration, complex           & Path amplitude $e^{iS/\hbar}$ & \textsc{con} \\
Sum over paths                & Feynman path integral & \textsc{con} \\
Interference of paths         & Quantum interference & \textsc{cnj} \\
No-arbitrage conversion       & Existence of an energy function & \textsc{thm} \\
Virtual token $\nu$           & Energy, as a Noether charge & \textsc{thm} \\
Holonomy of the energy form   & Failure of energy to be a state function & \textsc{thm} \\
Dissipation $\theta$          & Second law; free energy & \textsc{thm} \\
Ledger law $\sigma+\kappa=\sigma_0$ & Potential plus kinetic & \textsc{thm} \\
Erasure bounded by holonomy   & Landauer's principle & \textsc{thm} \\
Located authority             & Entanglement of space and time & \textsc{cnj} \\
$\nu$ generates reduction     & Hamiltonian, generator of evolution & \textsc{cnj} \\
Symmetry of the decoration    & Noether current & \textsc{cnj} \\
\bottomrule
\end{tabular}
\end{center}

\section{Running Example: The Rho Calculus}

Throughout this paper we use the Rho calculus~\cite{meredith2005rho} as a
running example because it exhibits, in compact form, many of the features of
interest: reflection (processes and names are mutually definable), parallel
composition, and a synchronization mechanism that serves as the prototype for
the causal graphs we study.

The grammar of the Rho calculus is:
\begin{align}
  P, Q &\;::=\; 0 \;\mid\; \mathtt{for}(y \leftarrow x)\,P \;\mid\;
               x!(Q) \;\mid\; P \mid Q \;\mid\; {*}x \\
  x, y &\;::=\; @P
\end{align}
where $y$ in $\mathtt{for}(y \leftarrow x)\,P$ is a bound variable.
Structural equivalence $\equiv$ is the smallest equivalence relation containing
$\alpha$-equivalence making $(P,\,|\,,0)$ a commutative monoid.
The single reduction rule is:
\[
  \mathtt{for}(y \leftarrow x)\,P \;\mid\; x!(Q)
  \;\rewrite\;
  P\{@Q / y\}
\]

\section{Organization}

Section~\ref{sec:gslt} defines GSLTs and their category.
Section~\ref{sec:reversibility} constructs the time-symmetric envelope $\GSLT^\dagger$.
Section~\ref{sec:causal} identifies the synchronization tree with a causal graph.
Section~\ref{sec:lts-hml} reviews the Milner--Sewell--Leifer construction and
the resulting context-decorated HML.
Section~\ref{sec:weighted} introduces the decoration and locates the
least-action principle in its tropical instance.
Section~\ref{sec:cost} asks what is conserved, and at what price.
Section~\ref{sec:extended-modal} extends the modal operators with full dynamical state.
Section~\ref{sec:quantum} lifts weights to complex amplitudes and connects to the
quantum Gillespie algorithm.
Section~\ref{sec:synthesis} collects the structures into a unified picture and
states the main conservation theorem.

% ===================================================================
